% Fuente: Auxiliar 4, Pregunta 2.
{\color{MidnightBlue}
\begin{enumerate}[label=\alph*)]

\item
Definimos las variables de decisión como:
\[
x_g := \text{producción de la turbina a gas [MW]},
\]
\[
x_v := \text{producción del bloque de vapor [MW]}.
\]

La energía producida por ambos componentes se remunera a un precio de
\(5\;[\$/\mathrm{MWh}]\). Además, los costos de operación son
\(4\;[\$/\mathrm{MWh}]\) para la turbina a gas y \(2\;[\$/\mathrm{MWh}]\)
para el bloque de vapor. Por lo tanto, el beneficio operativo es:
\[
5(x_g+x_v)-4x_g-2x_v.
\]

Simplificando:
\[
5x_g+5x_v-4x_g-2x_v
=
x_g+3x_v.
\]

Luego, el problema de optimización lineal queda dado por:
\[
\begin{aligned}
\max_{x_g,x_v} \quad & x_g+3x_v \\
\text{s.a.} \quad
& x_g \geq 4, \\
& x_g \leq 16, \\
& x_g+x_v \leq 20, \\
& x_v \leq 0.5x_g, \\
& x_g,x_v \geq 0.
\end{aligned}
\]

La primera restricción representa el mínimo de generación requerido para la
turbina a gas. La segunda corresponde al límite máximo de producción de la
turbina a gas. La tercera indica que la producción total de la central no puede
superar los \(20\;[\mathrm{MW}]\). Finalmente, la cuarta restricción representa
el acoplamiento físico del ciclo combinado, donde la producción del bloque de
vapor no puede exceder la mitad de la producción de la turbina a gas.

Para escribir el problema en forma estándar, primero transformamos la función
objetivo de maximización en una de minimización:
\[
\max \; x_g+3x_v
\quad \Longleftrightarrow \quad
\min \; -x_g-3x_v.
\]

Luego, incorporamos variables de holgura. Para la restricción
\(x_g \geq 4\), se introduce una variable \(s_1 \geq 0\):
\[
x_g - s_1 = 4.
\]

Para las restricciones de tipo menor o igual, se agregan variables de holgura
\(s_2,s_3,s_4 \geq 0\):
\[
x_g+s_2=16,
\]
\[
x_g+x_v+s_3=20,
\]
\[
-0.5x_g+x_v+s_4=0.
\]

Por lo tanto, la forma estándar es:
\[
\begin{aligned}
\min_{x_g,x_v,s_1,s_2,s_3,s_4} \quad
& -x_g-3x_v \\
\text{s.a.} \quad
& x_g-s_1 = 4, \\
& x_g+s_2 = 16, \\
& x_g+x_v+s_3 = 20, \\
& -0.5x_g+x_v+s_4 = 0, \\
& x_g,x_v,s_1,s_2,s_3,s_4 \geq 0.
\end{aligned}
\]

Equivalentemente, definiendo
\[
x =
\begin{pmatrix}
x_g\\
x_v\\
s_1\\
s_2\\
s_3\\
s_4
\end{pmatrix},
\]
el problema puede escribirse como:
\[
\begin{aligned}
\min_x \quad & c^\top x \\
\text{s.a.} \quad & Ax=b, \\
& x \geq 0,
\end{aligned}
\]
donde
\[
c =
\begin{pmatrix}
-1\\
-3\\
0\\
0\\
0\\
0
\end{pmatrix},
\]
\[
A =
\begin{pmatrix}
1 & 0 & -1 & 0 & 0 & 0 \\
1 & 0 & 0 & 1 & 0 & 0 \\
1 & 1 & 0 & 0 & 1 & 0 \\
-0.5 & 1 & 0 & 0 & 0 & 1
\end{pmatrix},
\qquad
b =
\begin{pmatrix}
4\\
16\\
20\\
0
\end{pmatrix}.
\]

\item
A partir del modelo de la parte anterior, se tiene el siguiente problema en las
variables originales:
\[
\begin{aligned}
\max_{x_g,x_v} \quad & x_g+3x_v \\
\text{s.a.} \quad
& x_g \geq 4, \\
& x_g \leq 16, \\
& x_g+x_v \leq 20, \\
& x_v \leq 0.5x_g, \\
& x_g,x_v \geq 0.
\end{aligned}
\]

La región factible está delimitada por las rectas
\[
x_g=4, 
\qquad
x_g=16,
\qquad
x_g+x_v=20,
\qquad
x_v=0.5x_g,
\qquad
x_v=0.
\]

\begin{figure}
    \centering
    \includegraphics[width=1\linewidth]{ejercicios/Capitulo-3/ejercicio_3_5_grap1.png}
    \caption{Método Gráfico P2.b}
    \label{ex:3:5:metodo-grafico}
\end{figure}


Los vértices de la región factible son:
\[
A=(4,0),
\]
\[
B=(16,0),
\]
\[
C=(16,4),
\]
\[
D=\left(\frac{40}{3},\frac{20}{3}\right),
\]
\[
E=(4,2).
\]

El punto \(D\) se obtiene de la intersección entre las restricciones activas
\[
x_g+x_v=20
\]
y
\[
x_v=0.5x_g.
\]
Resolviendo:
\[
x_g+0.5x_g=20,
\]
\[
1.5x_g=20,
\]
\[
x_g=\frac{40}{3}.
\]
Luego,
\[
x_v=0.5x_g
=
0.5\left(\frac{40}{3}\right)
=
\frac{20}{3}.
\]

Evaluando la función objetivo en cada vértice:
\[
z=x_g+3x_v,
\]
se obtiene:
\[
\begin{array}{c|c}
\text{Vértice} & z=x_g+3x_v \\ \hline
A=(4,0) & 4 \\
B=(16,0) & 16 \\
C=(16,4) & 28 \\
D=\left(\frac{40}{3},\frac{20}{3}\right) & \frac{100}{3} \\
E=(4,2) & 10
\end{array}
\]

Por lo tanto, el óptimo gráfico es
\[
(x_g^\star,x_v^\star)
=
\left(\frac{40}{3},\frac{20}{3}\right),
\]
con valor óptimo
\[
z^\star
=
\frac{100}{3}.
\]

En la forma estándar de minimización. Del análisis gráfico, en el óptimo se activan las restricciones
\[
x_g+x_v=20,
\qquad
-\frac{1}{2}x_g+x_v=0.
\]
Por lo tanto,
\[
s_3=0,
\qquad
s_4=0.
\]
Así, las variables no básicas son
\[
\mathcal{N}=\{s_3,s_4\}.
\]

En consecuencia, las variables básicas son
\[
\mathcal{B}=\{x_g,x_v,s_1,s_2\}.
\]

La matriz base \(B\) se obtiene tomando desde \(A\) las columnas asociadas a las variables básicas:
\[
B =
\begin{pmatrix}
\vert & \vert & \vert & \vert \\
A_{x_g} & A_{x_v} & A_{s_1} & A_{s_2} \\
\vert & \vert & \vert & \vert
\end{pmatrix}.
\]

Es decir,
\[
B =
\begin{pmatrix}
1 & 0 & -1 & 0 \\
1 & 0 & 0 & 1 \\
1 & 1 & 0 & 0 \\
-\frac{1}{2} & 1 & 0 & 0
\end{pmatrix}.
\]

Luego, la solución básica asociada a esta base se obtiene resolviendo
\[
x_B = B^{-1}b,
\]
donde
\[
x_B =
\begin{pmatrix}
x_g\\
x_v\\
s_1\\
s_2
\end{pmatrix},
\qquad
b =
\begin{pmatrix}
4\\
16\\
20\\
0
\end{pmatrix}.
\]

Por lo tanto,
\[
\begin{pmatrix}
x_g\\
x_v\\
s_1\\
s_2
\end{pmatrix}
=
B^{-1}b
=
\begin{pmatrix}
\frac{40}{3}\\
\frac{20}{3}\\
\frac{28}{3}\\
\frac{8}{3}
\end{pmatrix}.
\]

Como las variables no básicas son cero,
\[
s_3=0,
\qquad
s_4=0,
\]
la solución básica factible óptima completa es
\[
x^\star =
\begin{pmatrix}
x_g^\star\\
x_v^\star\\
s_1^\star\\
s_2^\star\\
s_3^\star\\
s_4^\star
\end{pmatrix}
=
\begin{pmatrix}
\frac{40}{3}\\
\frac{20}{3}\\
\frac{28}{3}\\
\frac{8}{3}\\
0\\
0
\end{pmatrix}.
\]

En dimensión mayor que dos no es práctico resolver el problema mediante el método
gráfico. En su lugar, se puede utilizar el método simplex, que recorre soluciones
básicas factibles del poliedro hasta encontrar una solución óptima.

La idea del método simplex es comenzar desde una solución básica factible, es decir, 
desde un vértice de la región factible. Luego, el algoritmo revisa si es posible 
mejorar el valor de la función objetivo moviéndose hacia un vértice vecino. Si existe 
una dirección que mejora la función objetivo, el método avanza hacia ese nuevo vértice.

Este procedimiento se repite hasta llegar a un vértice donde ya no es posible mejorar 
el valor de la función objetivo moviéndose a otro vértice vecino. En ese caso, dicho 
vértice corresponde a una solución óptima del problema.

\item
Ahora se agrega la condición operativa adicional:
\[
x_g+\frac{1}{2}x_v \geq 17.
\]

\begin{figure}
    \centering
    \includegraphics[width=1\linewidth]{ejercicios/Capitulo-3/ejercicio_3_5_mgp2c.png}
    \caption{Método Gráfico P2.c}
    \label{ex:3:5:restriccion-adicional}
\end{figure}

Por lo tanto, el problema actualizado queda dado por:
\[
\begin{aligned}
\max_{x_g,x_v} \quad & x_g+3x_v \\
\text{s.a.} \quad
& x_g \geq 4, \\
& x_g \leq 16, \\
& x_g+x_v \leq 20, \\
& x_v \leq \frac{1}{2}x_g, \\
& x_g+\frac{1}{2}x_v \geq 17, \\
& x_g,x_v \geq 0.
\end{aligned}
\]

Notemos que el óptimo anterior,
\[
\left(\frac{40}{3},\frac{20}{3}\right),
\]
ya no es factible, pues
\[
\frac{40}{3}
+
\frac{1}{2}\frac{20}{3}
=
\frac{50}{3}
<
17.
\]

La nueva restricción elimina parte de la región factible anterior. Los vértices
relevantes de la nueva región factible son:
\[
F=(16,2),
\]
\[
C=(16,4),
\]
\[
G=(14,6).
\]

Evaluamos la función objetivo original:
\[
z=x_g+3x_v.
\]

\[
\begin{array}{c|c}
\text{Vértice} & z=x_g+3x_v \\ \hline
F=(16,2) & 22 \\
C=(16,4) & 28 \\
G=(14,6) & 32
\end{array}
\]

Por lo tanto, el nuevo óptimo es
\[
(x_g^\star,x_v^\star)=(14,6),
\]
con valor óptimo
\[
z^\star=14+3(6)=32.
\]

Para actualizar la forma estándar, se agrega una variable de exceso
\(s_5\geq 0\) asociada a la nueva restricción:
\[
x_g+\frac{1}{2}x_v \geq 17.
\]

Por lo tanto,
\[
x_g+\frac{1}{2}x_v-s_5=17.
\]

Así, la forma estándar actualizada queda dada por:
\[
\begin{aligned}
\min \quad & -x_g-3x_v\\
\text{s.a.}\quad
& x_g-s_1=4,\\
& x_g+s_2=16,\\
& x_g+x_v+s_3=20,\\
& -\frac{1}{2}x_g+x_v+s_4=0,\\
& x_g+\frac{1}{2}x_v-s_5=17,\\
& x_g,x_v,s_1,s_2,s_3,s_4,s_5\geq 0.
\end{aligned}
\]


La solución básica factible óptima es
\[
x^\star=
\begin{pmatrix}
x_g\\
x_v\\
s_1\\
s_2\\
s_3\\
s_4\\
s_5
\end{pmatrix}
=
\begin{pmatrix}
14\\
6\\
10\\
2\\
0\\
1\\
0
\end{pmatrix}.
\]

Como en el óptimo se cumplen con igualdad las restricciones
\[
x_g+x_v=20,
\qquad
x_g+\frac{1}{2}x_v=17,
\]
se tiene que
\[
s_3=0,
\qquad
s_5=0.
\]

Luego, las variables no básicas son
\[
\mathcal{N}=\{s_3,s_5\},
\]
y la base óptima es
\[
\mathcal{B}=\{x_g,x_v,s_1,s_2,s_4\}.
\]

\item
Ahora el costo del bloque de vapor aumenta a
\[
c_v=6 \; [\$/\mathrm{MWh}].
\]

Por lo tanto, el beneficio operativo queda dado por
\[
5(x_g+x_v)-4x_g-6x_v=x_g-x_v.
\]

Luego, el nuevo problema es
\[
\max \; x_g-x_v
\]
sobre la misma región factible obtenida en la parte (c).

Los vértices factibles relevantes son
\[
F=(14,6), \qquad C=(16,4), \qquad E=(16,2).
\]

Evaluando la nueva función objetivo:
\[
\begin{array}{c|c}
\text{Vértice} & z=x_g-x_v \\ \hline
F=(14,6) & 8 \\
C=(16,4) & 12 \\
E=(16,2) & 14
\end{array}
\]

Por lo tanto, el nuevo óptimo es
\[
(x_g^\star,x_v^\star)=(16,2),
\]
con valor óptimo
\[
z^\star=14.
\]

En forma estándar, este punto se obtiene haciendo activas las restricciones
\[
x_g=16,
\qquad
x_g+\frac{1}{2}x_v=17.
\]

Por lo tanto,
\[
s_2=0,
\qquad
s_5=0.
\]

Además,
\[
s_1=x_g-4=12,
\]
\[
s_3=20-x_g-x_v=20-16-2=2,
\]
\[
s_4=\frac{1}{2}x_g-x_v=8-2=6.
\]

Así, la solución básica factible óptima es
\[
x^\star=
\begin{pmatrix}
x_g\\
x_v\\
s_1\\
s_2\\
s_3\\
s_4\\
s_5
\end{pmatrix}
=
\begin{pmatrix}
16\\
2\\
12\\
0\\
2\\
6\\
0
\end{pmatrix}.
\]

Comparando con la base óptima de la parte (c), se tiene que
\[
s_2 \text{ sale de la base},
\]
y
\[
s_3 \text{ entra a la base}.
\]

Por lo tanto, la nueva base óptima es
\[
\mathcal{B}=\{x_g,x_v,s_1,s_3,s_4\},
\]
mientras que las variables no básicas son
\[
\mathcal{N}=\{s_2,s_5\}.
\]

\item
Ahora se agrega la restricción térmica adicional:
\[
x_v \leq 4.5.
\]

Como se parte del problema de la parte (c), con \(c_v=2\), la función objetivo sigue siendo
\[
\max \; x_g+3x_v.
\]

El óptimo obtenido en la parte (c) era
\[
(x_g,x_v)=(14,6).
\]

Sin embargo, este punto ya no es factible, pues
\[
x_v=6>4.5.
\]

Por lo tanto, la base óptima anterior deja de ser factible.

La nueva restricción en forma estándar se escribe como
\[
x_v+s_6=4.5,
\qquad s_6\geq 0.
\]

El nuevo óptimo se obtiene en la intersección entre
\[
x_g+x_v=20
\]
y
\[
x_v=4.5.
\]

Así,
\[
x_g=20-4.5=15.5.
\]

Luego, el nuevo punto óptimo es
\[
(x_g^\star,x_v^\star)=(15.5,4.5).
\]

Evaluando la función objetivo:
\[
z^\star
=
x_g^\star+3x_v^\star
=
15.5+3(4.5)
=
29.
\]

En forma estándar, la solución queda dada por
\[
x^\star=
\begin{pmatrix}
x_g\\
x_v\\
s_1\\
s_2\\
s_3\\
s_4\\
s_5\\
s_6
\end{pmatrix}
=
\begin{pmatrix}
15.5\\
4.5\\
11.5\\
0.5\\
0\\
3.25\\
0.75\\
0
\end{pmatrix}.
\]

En efecto:
\[
s_1=x_g-4=11.5,
\]
\[
s_2=16-x_g=0.5,
\]
\[
s_3=20-x_g-x_v=0,
\]
\[
s_4=\frac{1}{2}x_g-x_v=3.25,
\]
\[
s_5=x_g+\frac{1}{2}x_v-17=0.75,
\]
\[
s_6=4.5-x_v=0.
\]

Por lo tanto, las restricciones activas son
\[
x_g+x_v=20,
\qquad
x_v=4.5.
\]

Así, las variables no básicas son
\[
\mathcal{N}=\{s_3,s_6\}.
\]

La nueva base óptima es
\[
\mathcal{B}
=
\{x_g,x_v,s_1,s_2,s_4,s_5\}.
\]

Por lo tanto, al agregar la restricción térmica, la solución óptima de la parte (c)
se vuelve infactible y el nuevo óptimo corresponde a
\[
(x_g^\star,x_v^\star)=(15.5,4.5),
\qquad
z^\star=29.
\]

\end{enumerate}
}

